\documentclass[12pt]{ctexart}
\usepackage{amsmath, amssymb}
\usepackage{geometry}
\usepackage{booktabs}
\usepackage{enumitem}
\geometry{margin=1in}
\newcommand{\colvec}[2]{\begin{pmatrix} #1 \\ #2 \end{pmatrix}}

\title{Answer Key}
\author{}
\date{}

\begin{document}
\maketitle

\section*{1\quad Vector translation review}
\begin{enumerate}[label=\textbf{\arabic*.}]
\item \((6,4)\).
\item \((3,6)\).
\end{enumerate}

\section*{2\quad Word embeddings}
\begin{enumerate}[label=\textbf{\arabic*.}]
\item \(\colvec{1}{3}\).
\item \(\colvec{1}{3}\).
\item The vector is the same in both cases. In this embedding it represents the \textbf{male-to-female / feminine} direction; it also maps \text{man} to \text{woman}.
\item Yes: \(\text{woman}-\text{man}=\colvec{3}{8}-\colvec{2}{5}=\colvec{1}{3}\), which matches Question 3.
\item Missing word: \textbf{princess}, at \((5,6)+\colvec{1}{3}=(6,9)\).
\end{enumerate}

\section*{3\quad Translating English to Chinese}
\begin{enumerate}[label=\textbf{\arabic*.}]
\item \(\mathbf{t}=\colvec{-15}{-10}\).
\item Completed table:
\begin{center}
\renewcommand{\arraystretch}{1.3}
\begin{tabular}{@{}llll@{}}
\toprule
English & English coords & Chinese & Chinese coords \\
\midrule
man     & (9,8)    & 男人   & (-6,-2) \\
woman   & (10,11)  & 女人   & (-5,1) \\
boy     & (7,8)    & 男孩   & (-8,-2) \\
girl    & (8,11)   & 女孩   & (-7,1) \\
king    & (8,6)    & 国王   & (-7,-4) \\
queen   & (9,9)    & 王妃   & (-6,-1) \\
prince  & (6,6)    & 王子   & (-9,-4) \\
princess& (7,9)    & 公主   & (-9,0) \\
\bottomrule
\end{tabular}
\end{center}
Missing grid labels: \((10,11)=\text{woman}\), \((9,9)=\text{queen}\), \((6,6)=\text{prince}\), \((7,9)=\text{princess}\), and \((-5,1)=\text{女人}\), \((-8,-2)=\text{男孩}\), \((-9,0)=\text{公主}\).
\item \(\colvec{-2}{0}\) represents the \textbf{adult/older to younger/child} direction (age/youth direction).
\item The pattern breaks for \text{princess}/公主 because 公主 is a separate lexicalised word, not a regular feminine form of 王子. Applying \(\mathbf{t}\) to \text{princess} gives \((-8,-1)\), but the actual 公主 is \((-9,0)\).
\end{enumerate}

\end{document}